\renewcommand*{\authorOfNextLines}{Helmut Quinz}
\begin{quote}
		\gqqEmph{Planung ersetzt Zufall durch Irrtum.}\\
	\small{Albert Einstein}
\end{quote}
Am Beginn eines Projekts steht dessen detaillierte Planung. 
Auch (nach Meinung des Autors vor allem) um im Projektverlauf zu erkennen ob es gefährdet ist zu scheitern.
Erwarten Sie nicht, dass das Projekt Ihrem Plan folgt.
Das Vorhandensein einer Planung verschafft den Beteiligten eine Ahnung über
\begin{itemize}
	\item benötigte materielle Ressourcen,
	\item erforderliche Tätigkeiten, 
	\item Verantwortlichkeiten (wer muss was machen),
	\item mögliche Projektrisiken \\und
	\item schafft die Möglichkeit \textbf{rechtzeitig} einem ungünstigen  Projektverlauf entgegenzusteuern.
\end{itemize}

Sofern nicht bereits ausreichende vorgenommen sind hier Projektumfang und Ziele, auch mögliche Ziele die explizit nicht verfolgt werden zu detaillieren.
Es ist ein vollständiger Projektplan mit
\begin{itemize}
	\item zu erfüllenden Aufgaben,
	\item erforderlichen materiellen Ressourcen (Finanzmittel, Werkzeuge, Werkstoffe, Bauteile, \dots),
	\item erforderliche immaterielle Ressourcen (Humankapital, Technologien, benötigtes Fachwissen), 
	\item feingliedrigen Zeitplan mit Dauer und Reihenfolge der einzelnen Arbeitsschritte. 
			Sie sollten hier auch zusätzliche Meilensteine einführen, an welche der Projektfortschritt gemessen werden kann
\end{itemize}
 	zu erstellen. \textbf{und}
\begin{itemize}
	\item Die Art der Überprüfung des Projektfortschritts festzulegen.
\end{itemize}
Berücksichtigen Sie bei der Zeitplanung Pufferzeiten für unvorhergesehene Ereignisse oder Verzögerungen!
Identifizieren Sie möglichst viele Risiken für das Projekt und bereiten sie Gegenstrategien vor.\\

Ein weiterer wichtiger Teil der Planung ist das Zuweisen erforderlicher Tätigkeiten an die einzelnen Team-Mitglieder.
Dazu gehört auch das Festlegen von Regeln für die interne Kommunikation \textbf{und} die mit externen Verantwortlichen (Betreuer)!\\

\paragraph{Wichtig:} Legen Sie hier \textbf{mit} Ihrem Betreuer Art und Häufigkeit der Fortschrittsrückmeldungen an den Betreuer und eventuell beteiligten externen Unternehmen fest!
\paragraph{Wichtig:} Achten Sie bei der Zeitplanung darauf dass sie spätestens ab den Herbstferien sich verstärkt um die allgemeinen schulischen Aufgaben kümmern werden müssen.
Es sind Schularbeiten zu schreiben, Tests zu bestehen und Prüfungen zu absolvieren. Zu diesem Zeitpunkt ist es hilfreich wenn ein Großteil der Projektarbeit \textbf{bereits erledigt} ist.
Andernfalls könnten der Abschluss des Jahrgangs und der, der \typeOfWork ~konkurrierende Ziele werden.